\chapter{Ricorsione}
\label{chap:ricorsione}

Le funzioni introdotte nella \cref{sec:fi-python-funzioni} possono
richiamare se stesse per risolvere un problema attraverso sottoproblemi
della stessa natura. Questa tecnica prende il nome di \emph{ricorsione}
e richiede di individuare un caso base e di ricavare la soluzione dai
risultati delle chiamate ricorsive.\footnote{Materiale
di riferimento: S. Canazza, \emph{Ricorsione}, file
\nolinkurl{2026_FI_03_canazza_Ricorsione.pdf}.}

\section{Funzioni ricorsive e terminazione}
\label{sec:fi-ricorsione-terminazione}

\begin{definition}[Ricorsione]
Una funzione è \emph{ricorsiva} se può invocare se stessa durante la
propria esecuzione. Un \emph{caso base} fornisce un risultato senza
ulteriori invocazioni ricorsive; un \emph{passo ricorsivo} esprime la
soluzione mediante uno o più sottoproblemi della stessa natura.
\end{definition}

Una chiamata ricorsiva avvia una nuova invocazione della funzione, con
parametri propri, mentre l'invocazione corrente rimane sospesa.
La presenza di un caso base non garantisce, da sola, la terminazione:
per ogni ingresso ammissibile, ciascuna successione di chiamate deve
raggiungere un caso base.

Un criterio sufficiente consiste nell'associare a ogni sottoproblema
una misura intera non negativa, strettamente decrescente a ogni
chiamata ricorsiva. Poiché non esistono successioni infinite strettamente
decrescenti di interi non negativi, tutte le chiamate raggiungono un
caso terminale, purché le singole operazioni non ricorsive terminino.
La terminazione, richiesta dalla \cref{def:fi-algoritmo}, non basta
a garantire la correttezza: occorre inoltre che i casi base e la
combinazione dei risultati forniscano la soluzione prevista.

\section{Il fattoriale: iterazione e ricorsione}
\label{sec:fi-fattoriale}

Per \(n\in\N\), il \emph{fattoriale} è definito dal prodotto
\(n!=1\cdot2\cdots n\), con \(0!=1\). Equivalentemente,
\begin{equation}
  n!=\begin{cases}
    1, & n=0,\\
    n(n-1)!, & n>0.
  \end{cases}
  \label{eq:fi-fattoriale}
\end{equation}

La definizione precedente si traduce nelle due implementazioni
seguenti, che calcolano lo stesso risultato per un intero non negativo:
la prima utilizza un ciclo (\cref{sec:fi-python-cicli}), la seconda
una chiamata ricorsiva.

\begin{unipdcode}[language=Python]
def fattoriale_iterativo(n):
    if n < 0:
        raise ValueError("n deve essere non negativo")
    risultato = 1
    for i in range(1, n + 1):
        risultato = risultato * i
    return risultato

def fattoriale(n):
    if n < 0:
        raise ValueError("n deve essere non negativo")
    if n == 0:
        return 1
    return n * fattoriale(n - 1)
\end{unipdcode}

Nella prima funzione, dopo l'iterazione di indice \(i\), la variabile
\texttt{risultato} vale \(i!\). Nella seconda, il caso base restituisce
\(0!=1\), mentre il passo ricorsivo applica l'\cref{eq:fi-fattoriale}.
La misura \(n\) diminuisce di uno a ogni chiamata e raggiunge zero.
La correttezza segue dalla stessa relazione: se la chiamata con \(n-1\)
restituisce \((n-1)!\), quella con \(n\) restituisce \(n!\).

\begin{unipdterminal}
>>> fattoriale(0), fattoriale(4)
(1, 24)
>>> fattoriale_iterativo(5), fattoriale(5)
(120, 120)
\end{unipdterminal}

Entrambe le funzioni presuppongono un argomento di tipo \texttt{int}.
Per \(n<0\) sollevano \texttt{ValueError}; senza tale controllo,
la versione ricorsiva continuerebbe a decrementare \(n\) senza
raggiungere il caso base.

\section{Invocazioni sospese e pila delle chiamate}
\label{sec:fi-pila-ricorsione}

Per descrivere l'esecuzione ricorsiva occorre considerare anche le
invocazioni non ancora terminate. Quando una funzione ne invoca un'altra,
la prima sospende la propria esecuzione fino al ritorno della seconda.
Ogni invocazione mantiene un
\emph{record di attivazione}, o \emph{frame}, contenente parametri,
variabili locali e informazioni necessarie a riprendere il calcolo.
Le invocazioni ancora attive costituiscono la \emph{pila delle chiamate}
(\emph{call stack}), organizzata secondo l'ordine LIFO
(\emph{Last In, First Out}), proprio dell'ADT pila della
\cref{sec:fi-adt-pila}.

\begin{example}[Traccia di \texttt{fattoriale(4)}]
\label{ex:fi-traccia-fattoriale}
Per valutare \texttt{fattoriale(4)}, le invocazioni con argomenti
\(4,3,2,1\) rimangono sospese fino al caso base \(n=0\), che restituisce
uno. Le chiamate riprendono quindi in ordine inverso:
\[
\begin{aligned}
  \texttt{fattoriale(0)}&=1,\\
  \texttt{fattoriale(1)}&=1\cdot1=1,\\
  \texttt{fattoriale(2)}&=2\cdot1=2,\\
  \texttt{fattoriale(3)}&=3\cdot2=6,\\
  \texttt{fattoriale(4)}&=4\cdot6=24.
\end{aligned}
\]
\end{example}

\begin{figure}[htbp]
  \centering
  \begin{tikzpicture}[unipd/diagram,x=1cm,y=1cm]
  \node at (0,0.65) {Chiamate};
  \node at (6,0.65) {Ritorni};
  \foreach \n/\y/\result in {4/0/24,3/-1/6,2/-2/2,1/-3/1,0/-4/1} {
    \node[unipd/node,minimum width=32mm] (c\n) at (0,\y)
      {\texttt{fattoriale(\n)}};
    \node[unipd/node,minimum width=32mm] (r\n) at (6,\y)
      {\(n=\n\): restituisce \(\result\)};
  }
  \foreach \n/\m in {4/3,3/2,2/1,1/0} {
    \draw[unipd/directed edge] (c\n) -- (c\m);
    \draw[unipd/directed edge] (r\m) -- (r\n);
  }
  \foreach \y/\text in {0/{\(4\cdot6=24\)},-1/{\(3\cdot2=6\)},-2/{\(2\cdot1=2\)},-3/{\(1\cdot1=1\)}} {
    \node at (3,\y) {\text};
  }
  \draw[unipd/directed edge] (c0) -- node[above] {Caso base: \(0!=1\)} (r0);
\end{tikzpicture}

  \caption{Invocazioni e ritorni nella valutazione di
  \texttt{fattoriale(4)}.}
  \label{fig:fi-fattoriale-chiamate}
\end{figure}

Nella \cref{fig:fi-fattoriale-chiamate}, al raggiungimento del caso base
sono presenti cinque frame. Le operazioni di moltiplicazione vengono
completate durante il ritorno delle chiamate: l'espressione
\texttt{n * fattoriale(n - 1)} richiede infatti il risultato della
chiamata interna.

Le due versioni del fattoriale effettuano \(n\) moltiplicazioni;
tuttavia, quella ricorsiva richiede \(n+1\) invocazioni e una
profondità di pila proporzionale a \(n\). La versione iterativa
utilizza un numero fisso
di variabili. Tale confronto considera il numero di operazioni e di
variabili; la dimensione degli interi prodotti cresce con \(n\).

\section{Ricorsione in coda ed eliminazione della ricorsione}
\label{sec:fi-ricorsione-coda}

La posizione della chiamata ricorsiva determina se rimangano
operazioni da eseguire dopo il suo ritorno. Una chiamata è \emph{in coda}
quando l'invocante non deve compiere ulteriori operazioni. La funzione
\texttt{fattoriale} della \cref{sec:fi-fattoriale} non è in coda,
poiché deve ancora moltiplicare il risultato per \texttt{n}.

Una versione in coda trasferisce il prodotto parziale in un parametro
\emph{accumulatore}:

\begin{unipdcode}[language=Python]
def fattoriale_coda(n, acc=1):
    if n < 0:
        raise ValueError("n deve essere non negativo")
    if n == 0:
        return acc
    return fattoriale_coda(n - 1, acc * n)
\end{unipdcode}

Il parametro \texttt{acc} ha valore iniziale \(1\). Per \(n=4\),
i parametri evolvono secondo la successione
\[
  (4,1)\longrightarrow(3,4)\longrightarrow(2,12)
  \longrightarrow(1,24)\longrightarrow(0,24).
\]

Poiché nessuna operazione rimane sospesa dopo la chiamata ricorsiva,
questa forma si trasforma direttamente in un ciclo:

\begin{unipdcode}[language=Python]
def fattoriale_ciclo(n, acc=1):
    if n < 0:
        raise ValueError("n deve essere non negativo")
    while n > 0:
        acc = acc * n
        n = n - 1
    return acc
\end{unipdcode}

È possibile eliminare anche una ricorsione non in coda, conservando
esplicitamente lo stato delle invocazioni sospese; in questo caso, però,
la trasformazione può richiedere una pila ausiliaria.

L'implementazione CPython non elimina automaticamente le chiamate
ricorsive in coda fra funzioni Python: ogni invocazione continua a
occupare un frame. Superato il limite di profondità delle chiamate,
l'interprete solleva \texttt{RecursionError}. Il limite dipende
dall'ambiente e non garantisce che un algoritmo matematicamente
terminante completi sempre l'esecuzione.\footnote{Documentazione Python,
\href{https://docs.python.org/3/library/sys.html\#sys.getrecursionlimit}{\emph{sys.getrecursionlimit}}.}

\section{Ricorsione multipla e numeri di Fibonacci}
\label{sec:fi-fibonacci}

Si parla di \emph{ricorsione multipla} quando una stessa invocazione
può richiamare ricorsivamente la funzione più di una volta. Un esempio
è il calcolo della successione di Fibonacci, definita per \(n\ge1\) da
\begin{equation}
  F_1=F_2=1,\qquad
  F_n=F_{n-1}+F_{n-2}\quad(n>2).
  \label{eq:fi-fibonacci}
\end{equation}

La traduzione diretta dell'\cref{eq:fi-fibonacci} richiede due casi
base e due chiamate ricorsive:

\begin{unipdcode}[language=Python]
def fib(n):
    if n < 1:
        raise ValueError("n deve essere almeno 1")
    if n <= 2:
        return 1
    return fib(n - 1) + fib(n - 2)
\end{unipdcode}

Per \(n>2\), entrambi gli argomenti delle chiamate sono positivi e
minori di \(n\); ogni ramo raggiunge quindi un caso base.
L'esecuzione delle due chiamate è sequenziale, non parallela.

\begin{figure}[htbp]
  \centering
  % Diagramma originale; licenza CC BY-SA 4.0 del materiale didattico.
\begin{tikzpicture}[unipd/diagram,x=1cm,y=1cm]
  \node[unipd/node] (a) at (0,0) {\texttt{fib(5)}};
  \node[unipd/node] (b) at (-2.4,-1.2) {\texttt{fib(4)}};
  \node[unipd/node] (c) at (2.4,-1.2) {\texttt{fib(3)}};
  \node[unipd/node] (d) at (-3.6,-2.4) {\texttt{fib(3)}};
  \node[unipd/node] (e) at (-1.2,-2.4) {\texttt{fib(2)}};
  \node[unipd/node] (f) at (1.2,-2.4) {\texttt{fib(2)}};
  \node[unipd/node] (g) at (3.6,-2.4) {\texttt{fib(1)}};
  \node[unipd/node] (h) at (-4.5,-3.6) {\texttt{fib(2)}};
  \node[unipd/node] (i) at (-2.7,-3.6) {\texttt{fib(1)}};
  \foreach \parent/\child in {a/b,a/c,b/d,b/e,c/f,c/g,d/h,d/i} {
    \draw[unipd/directed edge] (\parent) -- (\child);
  }
\end{tikzpicture}

  \caption{Albero delle nove invocazioni effettuate da
  \texttt{fib(5)}.}
  \label{fig:fi-fibonacci-albero}
\end{figure}

La \cref{fig:fi-fibonacci-albero} evidenzia il ricalcolo di uno
stesso termine: \texttt{fib(3)} viene eseguita due volte e
\texttt{fib(2)} tre. L'albero rappresenta tutte le invocazioni,
mentre la pila contiene soltanto quelle ancora attive.

Indicando con \(C_n\) il numero di invocazioni complessive, si ha
\[
  C_1=C_2=1,\qquad C_n=1+C_{n-1}+C_{n-2}\quad(n>2).
\]
Dalle due relazioni segue per induzione \(C_n=2F_n-1\).
Poiché \(F_{n+2}\ge2F_n\), il numero di chiamate cresce
esponenzialmente con \(n\), a causa del ricalcolo degli stessi
sottoproblemi. La profondità massima della pila cresce invece
linearmente. La stima dei costi temporali e spaziali è ripresa nella
\cref{sec:fi-complessita-fibonacci}.

Per evitare questi ricalcoli, un'implementazione iterativa conserva
soltanto due termini successivi:

\begin{unipdcode}[language=Python]
def fib_iterativo(n):
    if n < 1:
        raise ValueError("n deve essere almeno 1")
    precedente = 1
    corrente = 1
    for i in range(3, n + 1):
        nuovo = precedente + corrente
        precedente = corrente
        corrente = nuovo
    return corrente
\end{unipdcode}

Prima dell'iterazione di indice \(i\), le variabili
\texttt{precedente} e \texttt{corrente} valgono rispettivamente
\(F_{i-2}\) e \(F_{i-1}\). Si eseguono \(n-2\) somme per \(n>2\),
con un numero fisso di variabili, senza ricalcolare i termini.
Anche qui il confronto non considera il costo aritmetico degli interi
di dimensione crescente.

\clearpage

\begin{unipdterminal}
>>> fib(5), fib_iterativo(5)
(5, 5)
>>> fib(10), fib_iterativo(10)
(55, 55)
\end{unipdterminal}

La ricorsione multipla non comporta necessariamente ricalcoli:
l'ordinamento per fusione della \cref{sec:fi-ordinamento-fusione}
ne fornisce un esempio.

\section{Esercizi di riepilogo}

\begin{exercise}[Somma dei primi naturali]
\label{ex:fi-somma-naturali}
Scrivere una funzione ricorsiva che, per un intero \(n\ge0\), restituisca
\(1+2+\cdots+n\), con somma nulla per \(n=0\). Identificare il caso
base, il passo ricorsivo e una misura di terminazione.
\end{exercise}

\begin{proof}[Soluzione dell'esercizio]
Il caso base restituisce zero; il passo ricorsivo aggiunge \(n\) alla
somma dei primi \(n-1\) naturali.

\begin{unipdcode}[language=Python]
def somma_naturali(n):
    if n < 0:
        raise ValueError("n deve essere non negativo")
    if n == 0:
        return 0
    return n + somma_naturali(n - 1)
\end{unipdcode}

La misura \(n\) diminuisce di uno a ogni chiamata e raggiunge zero.
Per \(n=4\), la funzione restituisce \(10\).
\end{proof}

\begin{exercise}[Conto alla rovescia]
\label{ex:fi-conto-rovescia}
Scrivere una funzione ricorsiva che, per un intero \(n\ge0\), stampi
in ordine decrescente tutti gli interi da \(n\) a zero, uno per riga.
Distinguere i valori stampati da quello restituito.
\end{exercise}

\begin{proof}[Soluzione dell'esercizio]
Ogni chiamata stampa \(n\) e, se \(n>0\), richiama la funzione con
\(n-1\). Il caso base è \(n=0\).

\begin{unipdcode}[language=Python]
def conto_rovescia(n):
    if n < 0:
        raise ValueError("n deve essere non negativo")
    print(n)
    if n == 0:
        return
    return conto_rovescia(n - 1)
\end{unipdcode}

La misura \(n\) decresce fino a zero. Per \(n=3\), la funzione
stampa \(3,2,1,0\), uno per riga, e restituisce \texttt{None}.
L'invocazione ricorsiva è in coda secondo la
\cref{sec:fi-ricorsione-coda}.
\end{proof}
